\uuid{ii0g}
\exo7id{4514}
\titre{exo7 4514}
\auteur{quercia}
\organisation{exo7}
\datecreate{2010-03-14}
\isIndication{false}
\isCorrection{false}
\chapitre{Suite et série de fonctions}
\sousChapitre{Convergence simple, uniforme, normale}
\cpge{mpsi-mp2i.nombres-reels-suites.d, pcsi.nombres-reels-suites.c, ptsi.nombres-reels-suites.c, mpsi-mp2i.nombres-reels-suites.c, pcsi.nombres-reels-suites.b, ptsi.nombres-reels-suites.b, bcpst1.suites-reelles, mp-mpi.suites-series-fonctions.a, pc.suites-series-fonctions.a, psi.suites-series-fonctions.a}
\module{Analyse}
\niveau{L2}
\difficulte{}

\contenu{
\question{
On définit une suite de fonctions ${f_n} : {\R^{+*}} \to {\R^{+*}}$
par récurrence :
$\begin{cases} f_0(x) = x \cr
         f_{n+1}(x) = \frac 12\left({f_n(x) + \frac x{f_n(x)}}\right).\cr \end{cases}$

\'Etudier la convergence simple, puis uniforme des $f_n$
$\left({\text{considérer }
 g_n(x)= \frac {f_n(x) - \sqrt x}{f_n(x) + \sqrt x}}\right)$.
}
}
